\documentclass[11pt]{article}
\usepackage[a4paper,margin=25mm]{geometry}
\usepackage{fontspec}
\usepackage[hidelinks]{hyperref}
\setmainfont{Latin Modern Roman}
\pagestyle{plain}
\begin{document}
{\LARGE\bfseries ALIUS Bulletin - Guidelines for Commentary - 2022\par}
\vspace{1em}
\section*{Alius Bulletin}
\section*{*}
\section*{Guidelines For}
\section*{Commentary}
If you want to edit a Commentary for the ALIUS Bulletin, please read the following guidelines carefully.\par

Debate threads offer a platform to an author (or several authors) to critically discuss a recently published article or book dedicated to consciousness and its altered states.\par

Each debate thread consists of at least: (1) a commentary from the invited author(s) about the target piece. Each debate thread consists of at most: (1) a commentary from the invited author(s) about the target piece; (2) a reply to the commentary from the author(s) of the target piece; (3) a reply to the reply; (4) and a final reply.\par

Each commentary/reply will be first published as a separate blog article on the ALIUS Blog page. And then subsequently published as part of the ALIUS Bulletin.\par

Length of the commentary\par

Each Commentary (or reply to the commentary) must include the following:\par

-A title\par

-An abstract: 50 < 100 words\par

-Main text: 1000 < 2000 words\par

-References: no limit on the number of references; but the style used should be the 6th edition of the APA (if you use Zotero, this citation style can be downloaded at the following link: https://www.zotero.org/styles?q=id\%3Aapa).\par

-Addresses and affiliations of the authors for correspondence\par

-Font: Times New Roman; size: 12\par

Content of the commentary\par

Comments/replies are not expected to be formulated in rigorous academic writing style, but more as an invitation to clarify some positions. Thus, it is not intended to be an attack of the positions defended by the original author. They can fairly depart from the original approach of the target article in order to offer an alternative viewpoint. It also invites the original author for a reply clarifying her own understanding of this issue, and other scholars which are non-specialists to have a better understanding of some misconceptions appearing in the reception of some research. Each comment should be easily understandable to any scholar. For example, an anthropologist reading a comment conducted with a neuroscientist should easily understand most of the discussion, and conversely, a neuroscientist reading an interview conducted with an anthropologist should easily understand most of the discussion.\par

Modus operandi and requirements\par

1. Ask one of the editors of the Bulletin whether the researcher(s) you want to invite for a commentary fits well with the Bulletin’s editorial line.\par

2. Ask the researcher you want to invite to write a  commentary whether he/she agrees in principle to comment (when you do so, let him/her know about the general concept of the ALIUS Debates, the Bulletin, the aim of ALIUS, the modus operandi, the schedule, etc.).\par

3. Set up a deadline specifying when the commentary is expected to be received.\par

4. Once you receive the commentary, invite the original author(s) to write up a reply in the same format and explain him/her the general concept of the ALIUS Debates.\par

5. Once the original author(s) of the target piece is informed about the publication about the commentary, the commentary will be published on the Blog.\par

6. If the author(s) of the original target piece has agreed to write up a reply to the commentary, then set up a deadline (no more than a month), and when the reply is ready, send the reply to one of the editors for publication on the Blog.\par

7. Once the reply to the commentary has been received, invite the authors of the commentary to reply to the reply and set up a deadline.\par

8. Once you have received the reply to the reply, send it to one of the editors and ask the authors of the target piece whether they want to write a final reply. If so, set up a new deadline.\par

9. When all the back and forth commentaries and replies have been collected, you, as an editor of these commentaries and replies, should give a title to the whole series of the commentaries and replies. This title should start with the word “Debating…” (For example: “Debating the phenomenal vs. access consciousness dichotomy”).\par

10. Write up short bios of the commentator(s) and the original author(s) (for length and style of the bio, see the bios featured in the first issue of the Bulletin, in the section “Contributors”). You may use the official bio provided by the scholar’s website, the scholar’s university or publisher, etc.\par

11. Once they have received the commentaries/replies, the editors will forward it to a native English-speaking proofreader who will take care of checking grammar and linguistic correctness. Note that the job of proofreaders is not to correct typos. As indicated above, typos must be corrected before sending the file to the editors.\par

12. The editors will next send you the feedback provided by the proofreaders. You will be expected to integrate the changes suggested by the proofreaders.\par

13. When the final version is ready, send it to the commentator(s)/original author(s) to get their final agreement for publication.\par

14. Once you have the commentator(s)/original author(s) agreement for publication, send the final version of the Debate thread to one of the editors. The comment will then be ready to be published!\par

Editorial team\par

Daniel Friedman (danielarifriedman@gmail.com); Matthieu Koroma (matthieu.koroma@gmail.com) ; George Frejer (gyorgyf@live.com)\par

\end{document}
